% GENERATED FILE. Edit canonical lessons, then run npm run generate:books.
% canonical-source-hash: fnv1a64:2c08bfd131b8149a
% canonical-lessons: GE-C331-ueberzeugen, GE-C331-ueberzeugend, GE-C331-ueberzeugung, GE-C331-ueberreden, GE-C331-argumentieren

\chapter{Jemanden überzeugen}
\label{ch:ge-jemanden-uberzeugen}
\hlchaptermodality{331}

\begin{chapteropening}
\textbf{By the end of this chapter:} \emph{I can convince someone, talk someone into something, argue my case and say whether an answer is convincing.}

Complete the last lesson of chapter 331: I can convince someone, talk someone into something, argue my case and say whether an answer is convincing.
\end{chapteropening}

\section[überzeugen]{überzeugen — to convince}

\label{lesson:GE-C331-ueberzeugen}



\begin{quote}
Before the new one: say the German for the aspect, then the German for to advocate.
\end{quote}

\subsection*{You'll want to know: überzeugen}
\textbf{überzeugen} — \textquotedblleft{}to convince\textquotedblright{}.

\begin{quote}
Der Verkäufer hat mich überzeugt: Ich nehme das blaue Fahrrad.
\end{quote}

\begin{quote}
\emph{The salesman convinced me: I am taking the blue bike.}
\end{quote}

Inseparable: \emph{er überzeugt, er hat überzeugt}. \emph{Sich überzeugen} — to see for yourself that something is so.

\begin{cousinweb}
\emph{überzeugen} is built on \emph{zeugen} \textquotedblleft{}to testify\textquotedblright{}, from Middle High German \emph{ziugen} (EtymDB), with \emph{über-} \textquotedblleft{}over\textquotedblright{}. It is usually said to have begun in the law courts, for proving something against a person with witnesses.
\end{cousinweb}

\begin{grammarlens}[title={What you've built}]
One more word for telling, explaining and describing at B1.
\end{grammarlens}

\subsection*{Your turn}
\begin{itemize}
\raggedright
  \item \emph{Say it:} \emph{überzeugen}
  \item \emph{Say it:} \emph{überzeugen}, once more
  \item \emph{Say it:} a sentence of your own with \emph{überzeugen}
  \item \emph{Recall:} say the German for the aspect, then the German for to advocate, then say \emph{überzeugen} again
\end{itemize}

\begin{tcolorbox}[breakable,colback=teal!4,colframe=teal!35!black,title={Before you move on}]
What does überzeugen mean? (\textquotedblleft{}To convince\textquotedblright{}.) Say it once more.
\end{tcolorbox}
\section[überzeugend]{überzeugend — convincing}

\label{lesson:GE-C331-ueberzeugend}



\begin{quote}
Before the new one: say the German for to advocate, then the German for to convince.
\end{quote}

\subsection*{You'll want to know: überzeugend}
\textbf{überzeugend} — \textquotedblleft{}convincing\textquotedblright{}.

\begin{quote}
Seine Antwort war nicht sehr überzeugend.
\end{quote}

\begin{quote}
\emph{His answer was not very convincing.}
\end{quote}

The present participle of \emph{überzeugen}, here an adjective: \emph{ein überzeugendes Beispiel}.

\begin{cousinweb}
\emph{überzeugend} is the present participle of \emph{überzeugen}, with the ending \emph{-end}; the verb is built on \emph{zeugen} \textquotedblleft{}to testify\textquotedblright{}, from Middle High German \emph{ziugen} (EtymDB).
\end{cousinweb}

\begin{grammarlens}[title={What you've built}]
One more word for telling, explaining and describing at B1.
\end{grammarlens}

\subsection*{Your turn}
\begin{itemize}
\raggedright
  \item \emph{Say it:} \emph{überzeugend}
  \item \emph{Say it:} \emph{überzeugend}, once more
  \item \emph{Say it:} a sentence of your own with \emph{überzeugend}
  \item \emph{Recall:} say the German for to advocate, then the German for to convince, then say \emph{überzeugend} again
\end{itemize}

\begin{tcolorbox}[breakable,colback=teal!4,colframe=teal!35!black,title={Before you move on}]
What does überzeugend mean? (\textquotedblleft{}Convincing\textquotedblright{}.) Say it once more.
\end{tcolorbox}
\section[die Überzeugung]{die Überzeugung — the conviction, the firm belief}

\label{lesson:GE-C331-ueberzeugung}



\begin{quote}
Before the new one: say the German for to convince, then the German for convincing.
\end{quote}

\subsection*{You'll want to know: die Überzeugung}
\textbf{die Überzeugung} — \textquotedblleft{}the conviction, the firm belief\textquotedblright{}.

\begin{quote}
Er hat das aus Überzeugung gemacht, nicht für Geld.
\end{quote}

\begin{quote}
\emph{He did it out of conviction, not for money.}
\end{quote}

Plural: \emph{die Überzeugungen}. \emph{Ich bin der festen Überzeugung, dass …} — I firmly believe that …

\begin{cousinweb}
\emph{Überzeugung} is the noun of \emph{überzeugen} \textquotedblleft{}to convince\textquotedblright{}, formed with \emph{-ung}; the verb is built on \emph{zeugen} \textquotedblleft{}to testify\textquotedblright{}, from Middle High German \emph{ziugen} (EtymDB).
\end{cousinweb}

\begin{grammarlens}[title={What you've built}]
One more word for telling, explaining and describing at B1.
\end{grammarlens}

\subsection*{Your turn}
\begin{itemize}
\raggedright
  \item \emph{Say it:} \emph{die Überzeugung}
  \item \emph{Say it:} \emph{die Überzeugung}, once more
  \item \emph{Say it:} a sentence of your own with \emph{die Überzeugung}
  \item \emph{Recall:} say the German for to convince, then the German for convincing, then say \emph{die Überzeugung} again
\end{itemize}

\begin{tcolorbox}[breakable,colback=teal!4,colframe=teal!35!black,title={Before you move on}]
What does die Überzeugung mean? (\textquotedblleft{}The conviction, the firm belief\textquotedblright{}.) Say it once more.
\end{tcolorbox}
\section[überreden]{überreden — to talk (someone) into, to persuade}

\label{lesson:GE-C331-ueberreden}



\begin{quote}
Before the new one: say the German for convincing, then the German for the conviction.
\end{quote}

\subsection*{You'll want to know: überreden}
\textbf{überreden} — \textquotedblleft{}to talk (someone) into, to persuade\textquotedblright{}.

\begin{quote}
Meine Freundin hat mich überredet, mit ins Kino zu kommen.
\end{quote}

\begin{quote}
\emph{My friend talked me into coming along to the cinema.}
\end{quote}

Inseparable: \emph{sie überredet, sie hat überredet}. \emph{Überreden} works with pressure and many words; \emph{überzeugen} works with reasons.

\begin{cousinweb}
\emph{überreden} is built from \emph{über} \textquotedblleft{}over\textquotedblright{} and \emph{reden} \textquotedblleft{}to talk\textquotedblright{}, inherited from Middle High German \emph{reden} and Old High German \emph{redōn} \textquotedblleft{}to speak\textquotedblright{} (EtymDB): you talk someone over to your side.
\end{cousinweb}

\begin{grammarlens}[title={What you've built}]
One more word for telling, explaining and describing at B1.
\end{grammarlens}

\subsection*{Your turn}
\begin{itemize}
\raggedright
  \item \emph{Say it:} \emph{überreden}
  \item \emph{Say it:} \emph{überreden}, once more
  \item \emph{Say it:} a sentence of your own with \emph{überreden}
  \item \emph{Recall:} say the German for convincing, then the German for the conviction, then say \emph{überreden} again
\end{itemize}

\begin{tcolorbox}[breakable,colback=teal!4,colframe=teal!35!black,title={Before you move on}]
What does überreden mean? (\textquotedblleft{}To talk (someone) into, to persuade\textquotedblright{}.) Say it once more.
\end{tcolorbox}
\section[argumentieren]{argumentieren — to argue (a case), to give arguments}

\label{lesson:GE-C331-argumentieren}



\begin{quote}
Before the new one: say the German for the conviction, then the German for to talk.
\end{quote}

\subsection*{You'll want to know: argumentieren}
\textbf{argumentieren} — \textquotedblleft{}to argue (a case), to give arguments\textquotedblright{}.

\begin{quote}
Sie argumentiert gut, aber ich bleibe bei meiner Meinung.
\end{quote}

\begin{quote}
\emph{She argues well, but I am sticking to my opinion.}
\end{quote}

It means to reason, not to quarrel; a quarrel is \emph{der Streit}. Perfekt: \emph{hat argumentiert}, with no \emph{ge-}.

\begin{cousinweb}
\emph{argumentieren} is a learned borrowing from Latin \emph{argūmentārī} \textquotedblleft{}to bring forward proof, to prove\textquotedblright{} (Lewis \& Short), from \emph{argūmentum} \textquotedblleft{}proof\textquotedblright{}, with the German verb ending \emph{-ieren}.
\end{cousinweb}

\begin{grammarlens}[title={What you've built}]
One more word for telling, explaining and describing at B1.
\end{grammarlens}

\subsection*{Your turn}
\begin{itemize}
\raggedright
  \item \emph{Say it:} \emph{argumentieren}
  \item \emph{Say it:} \emph{argumentieren}, once more
  \item \emph{Say it:} a sentence of your own with \emph{argumentieren}
  \item \emph{Recall:} say the German for the conviction, then the German for to talk, then say \emph{argumentieren} again
\end{itemize}

\begin{tcolorbox}[breakable,colback=teal!4,colframe=teal!35!black,title={Before you move on}]
What does argumentieren mean? (\textquotedblleft{}To argue (a case), to give arguments\textquotedblright{}.) Say it once more.
\end{tcolorbox}
